% GENERATED FILE. Edit canonical lessons, then run npm run generate:books.
% canonical-source-hash: fnv1a64:35cb24a413405fea
% canonical-lessons: TE-C247-gillu, TE-C247-goku, TE-C247-gusagusaladu, TE-C247-oppincu, TE-C247-batimalu

\chapter{To pinch, To scratch, To whisper, To persuade, To plead}
\label{ch:te-to-pinch-to-scratch-to-whisper-to-persuade-to-plead}
\hlchaptermodality{247}

\begin{chapteropening}
\textbf{By the end of this chapter:} \emph{I can say to pinch, to scratch, to whisper, to persuade and to plead.}

Complete the last lesson of chapter 247: I can say to pinch, to scratch, to whisper, to persuade and to plead.
\end{chapteropening}

\section[\texorpdfstring{\te{గిల్లు} (gillu)}{గిల్లు (gillu)}]{\te{గిల్లు} (gillu) — to pinch}

\label{lesson:TE-C247-gillu}



\begin{quote}
Before the new one: say the Telugu for to decrease, then the Telugu for to refuse.
\end{quote}

\subsection*{You'll want to know: \te{గిల్లు}}
\textbf{\te{గిల్లు}} — \emph{gillu} — \textquotedblleft{}to pinch\textquotedblright{}.

\begin{grammarlens}[title={What you've built}]
One more word for this chapter.
\end{grammarlens}

\subsection*{Your turn}
\begin{itemize}
\raggedright
  \item \emph{Say it:} \emph{gillu}
  \item \emph{Say it:} \emph{gillu}, once more
  \item \emph{Say it:} say \emph{gillu}
  \item \emph{Recall:} say the Telugu for to decrease, then the Telugu for to refuse, then say \emph{gillu} again
\end{itemize}

\begin{tcolorbox}[breakable,colback=teal!4,colframe=teal!35!black,title={Before you move on}]
What does \te{గిల్లు} mean? (\textquotedblleft{}To pinch\textquotedblright{}.) Say it once more.
\end{tcolorbox}
\section[\texorpdfstring{\te{గోకు} (gōku)}{గోకు (gōku)}]{\te{గోకు} (gōku) — to scratch}

\label{lesson:TE-C247-goku}



\begin{quote}
Before the new one: say the Telugu for to refuse, then the Telugu for to pinch.
\end{quote}

\subsection*{You'll want to know: \te{గోకు}}
\textbf{\te{గోకు}} — \emph{gōku} — \textquotedblleft{}to scratch\textquotedblright{}.

\begin{grammarlens}[title={What you've built}]
One more word for this chapter.
\end{grammarlens}

\subsection*{Your turn}
\begin{itemize}
\raggedright
  \item \emph{Say it:} \emph{gōku}
  \item \emph{Say it:} \emph{gōku}, once more
  \item \emph{Say it:} say \emph{gōku}
  \item \emph{Recall:} say the Telugu for to refuse, then the Telugu for to pinch, then say \emph{gōku} again
\end{itemize}

\begin{tcolorbox}[breakable,colback=teal!4,colframe=teal!35!black,title={Before you move on}]
What does \te{గోకు} mean? (\textquotedblleft{}To scratch\textquotedblright{}.) Say it once more.
\end{tcolorbox}
\section[\texorpdfstring{\te{గుసగుసలాడు} (gusagusalāḍu)}{గుసగుసలాడు (gusagusalāḍu)}]{\te{గుసగుసలాడు} (gusagusalāḍu) — to whisper}

\label{lesson:TE-C247-gusagusaladu}



\begin{quote}
Before the new one: say the Telugu for to pinch, then the Telugu for to scratch.
\end{quote}

\subsection*{You'll want to know: \te{గుసగుసలాడు}}
\textbf{\te{గుసగుసలాడు}} — \emph{gusagusalāḍu} — \textquotedblleft{}to whisper\textquotedblright{}.

\begin{grammarlens}[title={What you've built}]
One more word for this chapter.
\end{grammarlens}

\subsection*{Your turn}
\begin{itemize}
\raggedright
  \item \emph{Say it:} \emph{gusagusalāḍu}
  \item \emph{Say it:} \emph{gusagusalāḍu}, once more
  \item \emph{Say it:} say \emph{gusagusalāḍu}
  \item \emph{Recall:} say the Telugu for to pinch, then the Telugu for to scratch, then say \emph{gusagusalāḍu} again
\end{itemize}

\begin{tcolorbox}[breakable,colback=teal!4,colframe=teal!35!black,title={Before you move on}]
What does \te{గుసగుసలాడు} mean? (\textquotedblleft{}To whisper\textquotedblright{}.) Say it once more.
\end{tcolorbox}
\section[\texorpdfstring{\te{ఒప్పించు} (oppiñcu)}{ఒప్పించు (oppiñcu)}]{\te{ఒప్పించు} (oppiñcu) — to persuade}

\label{lesson:TE-C247-oppincu}



\begin{quote}
Before the new one: say the Telugu for to scratch, then the Telugu for to whisper.
\end{quote}

\subsection*{You'll want to know: \te{ఒప్పించు}}
\textbf{\te{ఒప్పించు}} — \emph{oppiñcu} — \textquotedblleft{}to persuade\textquotedblright{}.

\begin{grammarlens}[title={What you've built}]
One more word for this chapter.
\end{grammarlens}

\subsection*{Your turn}
\begin{itemize}
\raggedright
  \item \emph{Say it:} \emph{oppiñcu}
  \item \emph{Say it:} \emph{oppiñcu}, once more
  \item \emph{Say it:} say \emph{oppiñcu}
  \item \emph{Recall:} say the Telugu for to scratch, then the Telugu for to whisper, then say \emph{oppiñcu} again
\end{itemize}

\begin{tcolorbox}[breakable,colback=teal!4,colframe=teal!35!black,title={Before you move on}]
What does \te{ఒప్పించు} mean? (\textquotedblleft{}To persuade\textquotedblright{}.) Say it once more.
\end{tcolorbox}
\section[\texorpdfstring{\te{బతిమాలు} (batimālu)}{బతిమాలు (batimālu)}]{\te{బతిమాలు} (batimālu) — to plead, to beg}

\label{lesson:TE-C247-batimalu}



\begin{quote}
Before the new one: say the Telugu for to whisper, then the Telugu for to persuade.
\end{quote}

\subsection*{You'll want to know: \te{బతిమాలు}}
\textbf{\te{బతిమాలు}} — \emph{batimālu} — \textquotedblleft{}to plead, to beg\textquotedblright{}.

\begin{grammarlens}[title={What you've built}]
One more word for this chapter.
\end{grammarlens}

\subsection*{Your turn}
\begin{itemize}
\raggedright
  \item \emph{Say it:} \emph{batimālu}
  \item \emph{Say it:} \emph{batimālu}, once more
  \item \emph{Say it:} say \emph{batimālu}
  \item \emph{Recall:} say the Telugu for to whisper, then the Telugu for to persuade, then say \emph{batimālu} again
\end{itemize}

\begin{tcolorbox}[breakable,colback=teal!4,colframe=teal!35!black,title={Before you move on}]
What does \te{బతిమాలు} mean? (\textquotedblleft{}To plead, to beg\textquotedblright{}.) Say it once more.
\end{tcolorbox}
